\documentclass[10pt,a4paper]{article}

\usepackage{fontspec}
\setmainfont{Cambria}
\setsansfont{Calibri}
\setmonofont{Consolas}[Scale=0.88]
\newfontfamily\thaifont{Leelawadee UI}[Scale=0.95]
\usepackage[a4paper,top=1.8cm,bottom=1.8cm,left=1.9cm,right=1.9cm,headheight=15pt]{geometry}
\usepackage{xcolor,graphicx,booktabs,longtable,array,caption,enumitem,amsmath,amssymb,fancyhdr,titlesec,etoolbox}
\usepackage[hidelinks]{hyperref}
\definecolor{ink}{HTML}{1F2A37}\definecolor{accent}{HTML}{1F4E79}\definecolor{muted}{HTML}{5B6676}
\definecolor{calloutbg}{HTML}{EEF3F9}\definecolor{calloutrule}{HTML}{1F4E79}
\hypersetup{colorlinks=true,linkcolor=accent,urlcolor=accent}
\setlength{\parindent}{0pt}\setlength{\parskip}{0.55em}
\linespread{1.06}
\setlist{itemsep=0.2em,topsep=0.3em,leftmargin=1.4em}
\captionsetup{font={small,color=muted},labelfont={bf,color=ink},justification=raggedright,singlelinecheck=false,skip=6pt}
\captionsetup[longtable]{position=top}
\newenvironment{callout}{\par\smallskip\noindent\begin{lrbox}{0}\begin{minipage}{\dimexpr\linewidth-1.6em\relax}}%
  {\end{minipage}\end{lrbox}\colorbox{calloutbg}{\hspace{0.3em}\usebox{0}\hspace{0.3em}}\par\smallskip}
\pretocmd{\tableofcontents}{\begingroup\hypersetup{linkcolor=ink}}{}{}
\apptocmd{\tableofcontents}{\endgroup}{}{}
\setlength{\emergencystretch}{2.5em}
\makeatletter
\@ifundefined{l@chapter}{}{\renewcommand*\l@section{\@dottedtocline{1}{1.5em}{2.6em}}
  \renewcommand*\l@subsection{\@dottedtocline{2}{4.1em}{3.4em}}}
\makeatother

\titleformat{\section}{\sffamily\large\bfseries\color{accent}}{\thesection}{0.6em}{}[{\color{accent}\titlerule[0.5pt]}]
\titleformat{\subsection}{\sffamily\normalsize\bfseries\color{ink}}{\thesubsection}{0.6em}{}
\titlespacing*{\section}{0pt}{10pt}{6pt}
\setlength{\parskip}{0.4em}
\pagestyle{fancy}\fancyhf{}
\fancyhead[L]{\sffamily\small\color{muted}AGENTRUN \textbullet{} interview quick reference}
\fancyhead[R]{\sffamily\small\color{muted}Alyssa L. \textbullet{} 12 November 2026}
\fancyfoot[C]{\sffamily\small\color{muted}\thepage}
\renewcommand{\headrulewidth}{0.4pt}
\begin{document}
{\sffamily\LARGE\bfseries\color{accent} AGENTRUN — interview quick reference\par}\vspace{4pt}All numbers are recomputed from raw outputs; full detail and sources are in the research dossier.

\subsection*{30-second introduction}
AGENTRUN is a multi-tool LLM agent runtime I built: a model calls tools such as search, weather and a calculator through MCP, and keeps a persistent memory across sessions. I asked whether that memory actually makes the agent more \emph{reliable}, not just more capable. On our institute's GPU cluster I ran controlled ablations: the same tasks with no memory, with memory, with a verifier, with provenance and with a rule-based staleness check, across models from 1.5 to 14 billion parameters. Memory roughly doubles task success. A verifier that cannot see where evidence came from makes things worse, and memory can stop an agent from saying "I don't know". Provenance, a simple staleness rule and two prompt rules fix most of that.

\subsection*{Two-minute explanation}
\begin{itemize}
\item \textbf{Origin.} QLoRA fine-tuning of Qwen2.5-7B on a 23-task benchmark on a Colab T4. One drop was partly a parser and client artefact; the later 23/23 came from an adapter whose extra training data was designed from failures on those same tasks. Re-run with a matched client: 23/23 five times, but no better than the base model on 55 fresh tasks (92.7\% vs 94.5\%).
\item \textbf{Question.} Memory had never been tested. Does it make a multi-tool agent more reliable, and which reliability mechanisms matter?
\item \textbf{Design.} Six scenarios (lost context, distractors, conflicts, stale records, unanswerable questions, tool failures). Five configurations adding one component each: A stateless, B memory, C + text verifier, D + provenance, E + staleness gate. Fictional entities, frozen tasks, a 77-check leakage audit, and scoring that separates success from harmful assertions.
\item \textbf{Result.} On 96 fresh tasks (Qwen2.5-7B and Qwen3-8B): 25\% $\rightarrow$ 57\% with memory, 52\% with a text verifier, 62\% with provenance, 75\% with the gate; stale records 4 $\rightarrow$ 25 of 32.
\item \textbf{Surprise.} With memory, Qwen3-8B stopped abstaining: it searched in a loop until the step limit. Two prompt rules fixed it and lifted the full system to 92\%.
\item \textbf{Capacity.} Inside Qwen2.5, the components helped most at 7B; 1.5B never uses memory, and 14B needs little beyond it.
\end{itemize}
\subsection*{Five strongest numerical findings}
\begingroup\small\setlength{\tabcolsep}{4pt}\renewcommand{\arraystretch}{1.18}
\begin{longtable}{>{\raggedright\arraybackslash}p{0.277\linewidth}>{\raggedright\arraybackslash}p{0.336\linewidth}>{\raggedright\arraybackslash}p{0.327\linewidth}}
\caption{The five findings to quote, with their denominators.}\\
\toprule
\textbf{Finding} & \textbf{Number} & \textbf{Denominator and test} \\
\midrule
\endfirsthead
\toprule
\textbf{Finding} & \textbf{Number} & \textbf{Denominator and test} \\
\midrule
\endhead
\bottomrule
\endlastfoot
Memory helps & 24.5\% $\rightarrow$ 56.8\% (+32.3 pp, 95\% CI +21.9 to +42.2) & 192 paired task blocks (96 fresh tasks × 2 models); McNemar p < 0.0001 \\
Text-only verifier hurts & −5.2 pp (original harness), −11.5 pp (revised) & same 192 pairs; p = 0.006 and p < 0.0001 \\
Provenance + staleness gate help & +10.4 pp then +13.0 pp; stale tasks 4 $\rightarrow$ 25 of 32 & same 192 pairs; p = 0.002 and p < 0.0001 \\
Memory suppresses abstention; a stopping rule fixes it & Qwen3-8B unanswerable questions: 15/16 (no memory), 1/16 (full system), 16/16 (with the rules); +15.8 pp overall & 16 tasks per cell; 480 paired runs per model, p < 0.0001 \\
Retrieval quality reaches the answers & 71.4\% $\rightarrow$ 41.7\% with 2,000 extra records (hash), 70.3\% with BM25 & 192 paired runs, Qwen2.5-7B; p < 0.0001 \\
\end{longtable}
\endgroup

\subsection*{Five major limitations}
\begin{enumerate}
\item \textbf{Synthetic memories.} Templated facts with clean entity–attribute keys, written by the harness; I designed both the method and the tasks.
\item \textbf{One trial per cell; Qwen-only local models.} The size sweep is one family; the hosted comparison was truncated by quota.
\item \textbf{Harmful answers are rare} (1–7 per 480 runs at 7–8B), so I cannot show a component reduces them.
\item \textbf{Incomplete HPC campaign.} 7 of 21 jobs never started: dense retrieval in the agent, the repair-only variant, the adapter on the memory benchmark, sampled repeats.
\item \textbf{No sealed-test result yet}, and an unexplained switch to Thai in the 14B model's reasoning.
\end{enumerate}
\subsection*{Ten difficult questions}
\textbf{Why did text-only verification hurt?} The same-model verifier saw the right record mixed with distractors and a "not found" from another tool, and called it unsupported; in campaign 1 it made 14 wrong rejections against 2 right ones. Self-correction without external feedback often fails (Huang et al., 2023).

\textbf{Then why keep a verifier?} It reliably rejects answers with no evidence: when Qwen2.5-7B invented 12500.00 without calling any tool, the verifier rejected it in all 12 runs.

\textbf{How did you prevent leakage?} Invented names; a check that none occurs in training data or other benchmarks; per-task memory stores verified after each run; ground truth never sent to the model, with every transcript scanned; freeze hashes the runner enforces; the revised harness built on dev tasks only. 77 checks pass.

\textbf{Why was the 23-task benchmark insufficient?} It was too small (one task = 4.3 points), run once per condition, used during development, and the V2 training data was designed from it. The fresh-task re-run shows the adapter's gain was benchmark-specific.

\textbf{Why does Qwen3 loop?} With a memory tool and no stopping rule, "not found" and near-miss records each invite another search. It never invented an answer; it ran out of steps (15 of 16 unanswerable questions in E).

\textbf{Why did the revised harness work?} One repair turn plus two rules (search memory before giving up; never repeat a call that answered, else abstain). For Qwen3-8B the repair turn never fired, so its +15.8 came from the rules. For Qwen2.5-7B both were active; the separating run did not complete.

\textbf{Why did false conflicts increase?} The models searched more and saw more old/new pairs; even with dates visible they called many of them conflicts (22, 22, 20 in B–D). Only the gate removed them (0 in E).

\textbf{What did the size sweep show?} An inverted U: 1.5B never calls memory; 3B gains from memory only; 7B gains from every component (best: 85/96); 14B gains mostly from memory, with calculator loops and Thai output as caveats.

\textbf{Does fine-tuning improve generalisation?} Not here: the V2 adapter scores 23/23 on the benchmark its extra data came from, but 92.7\% vs the base model's 94.5\% on 55 fresh tasks.

\textbf{What is novel?} Not the mechanisms (last-write-wins, verification, provenance all exist; see STALE, MemoryAgentBench). The contribution is a controlled, paired, replicated measurement of their marginal value inside a multi-tool agent across model sizes, including the negative result for verification without provenance.

\subsection*{Essential terminology}
\begingroup\small\setlength{\tabcolsep}{4pt}\renewcommand{\arraystretch}{1.18}
\begin{longtable}{>{\raggedright\arraybackslash}p{0.243\linewidth}>{\raggedright\arraybackslash}p{0.717\linewidth}}
\caption{Terms you may need to define on the spot.}\\
\toprule
\textbf{Term} & \textbf{One-line definition} \\
\midrule
\endfirsthead
\toprule
\textbf{Term} & \textbf{One-line definition} \\
\midrule
\endhead
\bottomrule
\endlastfoot
Configurations A–E & stateless; memory; + text-only verifier; + provenance; + deterministic staleness/conflict gate \\
Task block & one task under all five configurations by one model; the unit of every paired test \\
Harmful assertion & wrong value, stale value, unsupported claim or arbitrary pick between conflicting records \\
Safe miss & over-abstention or false conflict: not correct, but nothing false asserted \\
Provenance & record id, recording date and source attached to each memory result \\
Reliability gate & code that labels each record current, superseded, conflicting or unchallenged \\
Harness v1 / v2 & original prompt and loop / plus one repair turn and two stopping/search rules \\
McNemar test & paired test on the tasks where two configurations disagree \\
Task bootstrap & confidence interval from resampling tasks, keeping both models' results for a task together \\
QLoRA & LoRA adapter (here r = 16, 10.1M parameters) trained on a 4-bit quantised frozen base \\
Evaluation adaptation & designing training data from failures on the benchmark later used to score it \\
MCP & Model Context Protocol; in AGENTRUN it carries the tools, not the memory \\
\end{longtable}
\endgroup

\end{document}
